\section{用\texorpdfstring{$n$}{n}次抛物线弧近似代替其它曲线弧}
\markboth{84. 用$n$次抛物线弧近似代替其它曲线弧}{84. 用$n$次抛物线弧近似代替其它曲线弧}

在代数上，$n$ 次函数的一般形式是
\[
  y = a_0 + a_1 x + a_2 x^2 + \cdots + a_n x^n. \quad (a_n \neq 0)
\]
$n$ 次函数的图象，叫做 $n$ 次抛物线。

在很多与曲线有关的计算问题中，$n$ 次抛物线比其它曲线显得方便。例如，在描点作图时，需要根据 $x$ 的值计算 $y$ 的对应值；求切线斜率时需要根据 $x$ 的值计算 $y'$ 的对应值；求曲边梯形面积时要计算积分 $\int_a^b y \dd x$，等等。在上述这些计算中，$n$ 次抛物线都比较方便，所以在具体计算时，往往在精确度允许范围内，选择一段 $n$ 次抛物线弧来近似代替已知曲线弧。

由 $n$ 次函数连续求导，顺次得到
\begin{align*}
  y' &= a_1 + 2a_2 x + 3a_3 x^2 + \cdots + n a_n x^{n-1}, \\
  y'' &= 2a_2 + 6a_3 x + \cdots + n(n-1) a_n x^{n-2}, \\
  y''' &= 6a_3 + \cdots + n(n-1)(n-2) a_n x^{n-3}, \\
  &\cdots\cdots\cdots\cdots\cdots\cdots\cdots\cdots\cdots\cdots \\
  y^{(n)} &= n(n-1)(n-2)\cdots 3 \cdot 2 \cdot 1 \cdot a_n.
\end{align*}
当 $x = 0$ 时得到
\begin{align*}
  &y(0) = a_0,\quad y'(0) = a_1,\quad y''(0) = 2a_2 = 2! a_2, \\
  &y'''(0) = 6a_3 = 3! a_3,\quad \cdots, \\
  &y^{(n)}(0) = n! a_n.
\end{align*}
因而各个系数可用 $x = 0$ 时的函数值和各阶导数值表示如下：
\begin{align*}
  &a_0 = y(0),\quad a_1 = y'(0),\quad a_2 = \frac{y''(0)}{2!}, \\
  &a_3 = \frac{y'''(0)}{3!},\quad \cdots\cdots, \\
  &a_n = \frac{y^{(n)}(0)}{n!}.
\end{align*}

一般地，任意给定一条曲线，把它的方程写成形式
\[
  y = f(x),
\]
求出当 $x = 0$ 时的函数值 $f(0)$ 和各阶导数值 $f'(0)$、$f''(0)$、$\cdots$、$f^{(n)}(0)$，然后按上面的公式求出系数 $a_0$、$a_1$、$a_2$、$\cdots$、$a_n$，这样就得到一个 $n$ 次多项式
\begin{align*}
  y &= f(0) + \frac{f'(0)}{1!}x + \frac{f''(0)}{2!}x^2 + \frac{f'''(0)}{3!}x^3 + \cdots \\
    &\quad + \frac{f^{(n)}(0)}{n!}x^n,
\end{align*}
如果随着 $n$ 的增大，$\frac{f^{(n)}(0)}{n!}x^n \to 0$，那么就可以用这个 $n$ 次多项式近似地代替已知函数 $y = f(x)$。反映在图形上，就是用 $x = 0$ 附近的一段 $n$ 次抛物线的弧近似代替已知曲线弧。$n$ 越大，近似的程度就越好。

例如，如果已知函数是
\[
  y = e^x,
\]
那么 $y$ 的各阶导数仍为 $e^x$，因而当 $x = 0$ 时的 $y$ 值和各阶导数值都是 $1$，按上述方法作出的 $n$ 次函数是
\[
  y = 1 + \frac{x}{1!} + \frac{x^2}{2!} + \frac{x^3}{3!} + \cdots + \frac{x^n}{n!}.
\]
由于 $n \to \infty$ 时，$\frac{x^n}{n!} \to 0$，所以这个 $n$ 次多项式可以近似地代替已知函数 $y = e^x$。从图形上看，就是在 $x = 0$ 附近上述方程所表示的 $n$ 次抛物线可以近似代替指数曲线 $y = e^x$。% 原书此处漏印“次”字
